\chapter{Áreas de Atuação}

O profissional poderá atuar em organizações públicas ou privadas que desenvolvam ou utilizem Inteligência Artificial em seus processos.

Entre os principais segmentos destacam-se:

\textbf{Tecnologia da Informação}
\begin{itemize}
    \item Empresas de desenvolvimento de software.
    \item Software Houses.
    \item Empresas SaaS.
    \item Consultorias.
\end{itemize}

\textbf{Ciência de Dados}
\begin{itemize}
    \item Empresas de Analytics.
    \item Business Intelligence.
    \item Big Data.
    \item Data Engineering.
\end{itemize}

\textbf{Inteligência Artificial}
\begin{itemize}
    \item Empresas especializadas em IA.
    \item Startups de IA.
    \item Centros de Pesquisa.
    \item Laboratórios de Inovação.
\end{itemize}

\textbf{Indústria 4.0}
\begin{itemize}
    \item Automação Industrial.
    \item Manufatura Inteligente.
    \item Sistemas Ciberfísicos.
    \item Gêmeos Digitais.
\end{itemize}

\textbf{Saúde}
\begin{itemize}
    \item Diagnóstico Assistido por IA.
    \item Processamento de Imagens Médicas.
    \item Apoio à Decisão Clínica.
\end{itemize}

\textbf{Agronegócio}
\begin{itemize}
    \item Agricultura de Precisão.
    \item Monitoramento Inteligente.
    \item Análise Preditiva.
\end{itemize}

\textbf{Mercado Financeiro}
\begin{itemize}
    \item Análise de Risco.
    \item Detecção de Fraudes.
    \item Sistemas Inteligentes de Crédito.
\end{itemize}

\textbf{Logística}
\begin{itemize}
    \item Otimização de Rotas.
    \item Previsão de Demanda.
    \item Cadeia de Suprimentos Inteligente.
\end{itemize}

\textbf{Educação}
\begin{itemize}
    \item Plataformas Inteligentes.
    \item Tutores Virtuais.
    \item Sistemas Adaptativos.
\end{itemize}

\textbf{Setor Público}
\begin{itemize}
    \item Cidades Inteligentes.
    \item Governo Digital.
    \item Automação de Processos.
\end{itemize}
